\documentclass{itmo}

\setlist[itemize]{label=---}
\setlist[enumerate]{label=\arabic*)}

\usepackage{dcolumn}
\newcolumntype{d}[1]{D{.}{.}{#1}}

\begin{document}

\makereporttitle{}
	{лабораторной работе №4}
	{Структуры и алгоритмы в базах данных и распределенных системах}
	{Реализация конкурентной хэш-таблицы}
	{}
	{Постнов~С.~А./P4135}
	{Платонов~А.~В./доцент, Портнов~П.~В./ассистент}

\maketableofcontents

\chapter{Теоретическая часть}

\section{Конкурентные хэш-таблицы}

Конкурентная хэш-таблица~---~структура данных, допускающая одновременный доступ из нескольких потоков выполнения без нарушения целостности данных.
В отличие от обычной хэш-таблицы, где для синхронизации можно использовать единственную блокировку (coarse-grained locking), конкурентные реализации стремятся минимизировать конкуренцию за блокировки, чтобы обеспечить масштабируемость при увеличении числа потоков.

Основные подходы к реализации конкурентных хэш-таблиц:
\begin{itemize}
	\item \textit{Striped locking}~---~множество независимых блокировок, каждая из которых защищает свою полосу (группу бакетов).
	Количество блокировок фиксировано и не зависит от размера таблицы, что упрощает реализацию, но может приводить к ложной конкуренции при ресайзе;
	\item \textit{Lock striping с переносимыми блокировками}~---~аналог striped locking, но при ресайзе блокировки переносятся вместе с бакетами, что позволяет избежать глобальной блокировки;
	\item \textit{Lock-free реализации}~---~используют атомарные операции (CAS) для обновления бакетов без блокировок.
	Обеспечивают максимальную масштабируемость, но сложны в реализации, особенно при ресайзе;
	\item \texttt{sync.Map} (Go)~---~встроенная конкурентная карта Go, использующая двухуровневую схему: read-only кэш для быстрого чтения и dirty-карта для записи с амортизированным переносом.
\end{itemize}

\section{Striped locking}

В данной работе используется подход \textit{striped locking}: каждый бакет таблицы защищён собственным \texttt{sync.RWMutex}.
Это позволяет нескольким потокам одновременно читать из разных бакетов, а запись блокирует только один бакет.
Для ресайза используется отдельная блокировка \texttt{resizeMu}: при добавлении элемента поток захватывает \texttt{resizeMu.RLock()}, что позволяет нескольким потокам одновременно выполнять \texttt{Put}, а при необходимости ресайза~---~\texttt{resizeMu.Lock()}, блокируя все операции добавления на время перестройки таблицы.

Текущая таблица хранится в \texttt{atomic.Value}, что обеспечивает lock-free доступ к указателю на таблицу при чтении.
При ресайзе создаётся новая таблица увеличенного размера, после чего указатель атомарно заменяется.

\section{Хэш-функция FNV-1a}

FNV-1a (Fowler--Noll--Vo)~---~некриптографическая хэш-функция, хорошо подходящая для хэш-таблиц благодаря быстрому вычислению и равномерному распределению.
Алгоритм: для каждого байта входных данных выполняется операция XOR с текущим хэшем, затем результат умножается на простое число (prime).
Начальное значение хэша (offset basis) и множитель выбраны эмпирически для минимизации коллизий.

\chapter{Практическая часть}

\section{Архитектура реализации}

Основной код расположен в файле \texttt{concurrent/map.go}.
Структура данных состоит из следующих компонентов:
\begin{enumerate}
	\item \texttt{Map}~---~основная структура, содержащая:
	\begin{itemize}
		\item \texttt{resizeMu sync.RWMutex}~---~блокировка ресайза (чтение при \texttt{Put}, запись при \texttt{maybeResize});
		\item \texttt{table atomic.Value}~---~указатель на текущую таблицу (lock-free чтение);
		\item \texttt{size int64}~---~счётчик элементов (атомарный);
		\item \texttt{maxLoadFactor float64}~---~порог загрузки (по умолчанию 0.75).
	\end{itemize}
	\item \texttt{table}~---~контейнер для массива бакетов, хранящийся в \texttt{atomic.Value};
	\item \texttt{bucket}~---~бакет с \texttt{sync.RWMutex} и слайсом \texttt{[]entry};
	\item \texttt{entry}~---~пара ключ-значение (\texttt{key string}, \texttt{value string}).
\end{enumerate}

\section{Базовая реализация}

Файл \texttt{concurrent/map\_baseline.go} содержит простую обёртку над \texttt{map[string]string} без синхронизации, используемую как базовый вариант для сравнения в однопоточных бенчмарках.
Для многопоточного сравнения используется стандартная \texttt{sync.Map} (обёртка \texttt{MapSync} в \texttt{benchmark\_test.go}).

\chapter{Исследовательская часть}

\section{Аппаратные характеристики}

Все замеры проводились на ноутбуке со следующей конфигурацией:
\begin{itemize}
	\item операционная система: \texttt{macOS 26.3.1} (arm64);
	\item процессор: \texttt{Apple M3 Pro}, 11 ядер;
	\item оперативная память: \texttt{18 ГБ};
	\item версия Go: \texttt{go1.25}.
\end{itemize}
Во время замеров не выполнялись ресурсоёмкие фоновые задачи, а ноутбук был подключён к сети электропитания для обеспечения максимальной производительности.

\section{Методика замеров}

Для замеров использованы встроенные benchmark-тесты Go (\texttt{testing.B} с \texttt{b.Loop()}).
Измерялись две операции: \texttt{put} и \texttt{get} в двух режимах:
\begin{itemize}
	\item \textit{single}~---~однопоточное выполнение (сравнение с \texttt{MapBaseline});
	\item \textit{multi}~---~многопоточное выполнение с числом потоков, равным \texttt{GOMAXPROCS} (сравнение с \texttt{sync.Map}).
\end{itemize}

Для каждого размера набора данных ($N = 1000, 2500, 5000, 7500, 10000, 25000, 50000, 100000$) фиксировались:
\begin{itemize}
	\item средняя задержка на операцию (нс/оп);
	\item средняя пропускная способность (оп/с);
	\item 95\% доверительный интервал (CI95) для задержки и пропускной способности;
	\item потребление кучи на операцию (байт/оп).
\end{itemize}

Доверительный интервал вычисляется как $CI_{95} = 1.96 \cdot \sigma / \sqrt{n}$, где $\sigma$~---~стандартное отклонение выборки, $n$~---~число итераций.
Потребление памяти измеряется через \texttt{runtime.ReadMemStats} как разница \texttt{HeapAlloc} до и после выполнения пакета операций.

Для профилирования CPU используется отдельная программа \texttt{cmd/profile/main.go}, которая выполняет заданное число раундов операций с включённым \texttt{pprof.StartCPUProfile}.
Профили визуализируются как ориентированные графы вызовов с помощью \texttt{go tool pprof -dot} и \texttt{dot -Tpdf}.

\section{Результаты для одного потока}

Таблица~\ref{tab:single-latency} содержит результаты замеров задержки в однопоточном режиме.

\begin{table}[H]
	\centering
	\caption{Задержка операций в однопоточном режиме (нс/оп)}
	\label{tab:single-latency}
	\begin{tabularx}{\textwidth}{|r|X|X|X|X|}
		\toprule
		$N$ & вставка (concurrent) & вставка (baseline) & поиск (concurrent) & поиск (baseline) \\
		\midrule
		1000 & $97.10 \pm 0.16$ & $45.31 \pm 0.04$ & $14.31 \pm 0.01$ & $10.14 \pm 0.00$ \\
		2500 & $89.68 \pm 0.08$ & $42.45 \pm 0.04$ & $14.71 \pm 0.04$ & $10.90 \pm 0.01$ \\
		5000 & $88.69 \pm 0.08$ & $39.54 \pm 0.03$ & $14.04 \pm 0.01$ & $11.10 \pm 0.01$ \\
		7500 & $107.09 \pm 0.11$ & $45.48 \pm 0.05$ & $14.46 \pm 0.01$ & $10.56 \pm 0.01$ \\
		10000 & $94.41 \pm 0.69$ & $38.53 \pm 0.03$ & $14.34 \pm 0.01$ & $10.63 \pm 0.01$ \\
		25000 & $216.08 \pm 1.44$ & $55.14 \pm 0.12$ & $16.68 \pm 0.10$ & $11.51 \pm 0.01$ \\
		50000 & $189.63 \pm 1.31$ & $52.66 \pm 0.45$ & $19.50 \pm 0.17$ & $11.59 \pm 0.01$ \\
		100000 & $203.46 \pm 3.42$ & $58.76 \pm 0.39$ & $33.55 \pm 0.24$ & $12.00 \pm 0.04$ \\
		\bottomrule
	\end{tabularx}
\end{table}

Таблица~\ref{tab:single-throughput} содержит результаты замеров пропускной способности в однопоточном режиме.

\begin{table}[H]
	\centering
	\caption{Пропускная способность в однопоточном режиме (млн оп/с)}
	\label{tab:single-throughput}
	\begin{tabularx}{\textwidth}{|r|X|X|X|X|}
		\toprule
		$N$ & вставка (concurrent) & вставка (baseline) & поиск (concurrent) & поиск (baseline) \\
		\midrule
		1000 & $10.36 \pm 0.01$ & $22.19 \pm 0.02$ & $70.27 \pm 0.03$ & $99.10 \pm 0.04$ \\
		2500 & $11.16 \pm 0.01$ & $23.60 \pm 0.02$ & $68.74 \pm 0.07$ & $92.14 \pm 0.06$ \\
		5000 & $11.28 \pm 0.01$ & $25.31 \pm 0.02$ & $71.47 \pm 0.06$ & $90.63 \pm 0.08$ \\
		7500 & $9.34 \pm 0.01$ & $22.01 \pm 0.03$ & $69.33 \pm 0.06$ & $94.96 \pm 0.08$ \\
		10000 & $10.67 \pm 0.04$ & $25.97 \pm 0.02$ & $69.84 \pm 0.06$ & $94.27 \pm 0.08$ \\
		25000 & $4.64 \pm 0.03$ & $18.15 \pm 0.04$ & $61.13 \pm 0.29$ & $86.95 \pm 0.07$ \\
		50000 & $5.28 \pm 0.03$ & $19.14 \pm 0.15$ & $52.12 \pm 0.33$ & $86.30 \pm 0.07$ \\
		100000 & $4.93 \pm 0.08$ & $17.06 \pm 0.11$ & $29.92 \pm 0.17$ & $83.48 \pm 0.16$ \\
		\bottomrule
	\end{tabularx}
\end{table}

Графики результатов представлены на рисунках~\ref{img:latency_single} и~\ref{img:throughput_single}.

\includeimage
	{latency_single}
	{f}
	{H}
	{\textwidth}
	{Задержка операций в однопоточном режиме}

\includeimage
	{throughput_single}
	{f}
	{H}
	{\textwidth}
	{Пропускная способность в однопоточном режиме}

В однопоточном режиме реализация на основе \texttt{sync.RWMutex} ожидаемо медленнее \texttt{map[string]string} примерно в 2--4 раза для \texttt{put} и в 1.4--2.8 раза для \texttt{get}, что объясняется накладными расходами на захват/освобождение мьютексов и атомарные операции.
Задержка \texttt{put} растёт с увеличением $N$ из-за ресайзов и увеличения длины цепочек в бакетах, тогда как \texttt{get} остаётся стабильным до $N \approx 25000$, после чего начинает расти из-за ухудшения локальности кэша.

\section{Результаты для нескольких потоков}

Таблица~\ref{tab:multi-latency} содержит результаты замеров задержки в многопоточном режиме.

\begin{table}[H]
	\centering
	\caption{Задержка операций в многопоточном режиме (нс/оп)}
	\label{tab:multi-latency}
	\begin{tabularx}{\textwidth}{|r|X|X|X|X|}
		\toprule
		$N$ & вставка (concurrent) & вставка (syncmap) & поиск (concurrent) & поиск (syncmap) \\
		\midrule
		1000 & $144.76 \pm 0.22$ & $65.75 \pm 0.12$ & $38.63 \pm 0.06$ & $37.17 \pm 0.05$ \\
		2500 & $151.08 \pm 0.22$ & $61.23 \pm 0.15$ & $42.95 \pm 0.15$ & $41.02 \pm 0.15$ \\
		5000 & $162.02 \pm 0.40$ & $61.41 \pm 0.10$ & $48.39 \pm 0.22$ & $44.66 \pm 0.19$ \\
		7500 & $191.48 \pm 0.59$ & $62.33 \pm 0.10$ & $48.64 \pm 0.27$ & $47.84 \pm 0.23$ \\
		10000 & $180.94 \pm 0.59$ & $62.97 \pm 0.56$ & $54.01 \pm 0.22$ & $51.95 \pm 0.12$ \\
		25000 & $286.90 \pm 2.25$ & $80.41 \pm 0.28$ & $57.45 \pm 0.10$ & $51.04 \pm 0.09$ \\
		50000 & $305.69 \pm 1.49$ & $77.61 \pm 0.51$ & $58.24 \pm 0.11$ & $51.32 \pm 0.13$ \\
		100000 & $309.06 \pm 1.91$ & $83.11 \pm 0.77$ & $60.03 \pm 0.10$ & $53.83 \pm 0.24$ \\
		\bottomrule
	\end{tabularx}
\end{table}

\begin{table}[H]
	\centering
	\caption{Пропускная способность в многопоточном режиме (млн оп/с)}
	\label{tab:multi-throughput}
	\begin{tabularx}{\textwidth}{|r|X|X|X|X|}
		\toprule
		$N$ & вставка (concurrent) & вставка (syncmap) & поиск (concurrent) & поиск (syncmap) \\
		\midrule
		1000 & $6.94 \pm 0.01$ & $15.45 \pm 0.03$ & $26.44 \pm 0.04$ & $27.27 \pm 0.03$ \\
		2500 & $6.63 \pm 0.01$ & $16.52 \pm 0.04$ & $24.07 \pm 0.09$ & $25.27 \pm 0.10$ \\
		5000 & $6.19 \pm 0.02$ & $16.33 \pm 0.03$ & $21.30 \pm 0.12$ & $23.06 \pm 0.12$ \\
		7500 & $5.23 \pm 0.02$ & $16.07 \pm 0.03$ & $21.21 \pm 0.15$ & $21.22 \pm 0.10$ \\
		10000 & $5.54 \pm 0.02$ & $16.00 \pm 0.04$ & $18.74 \pm 0.10$ & $19.33 \pm 0.06$ \\
		25000 & $3.49 \pm 0.03$ & $12.46 \pm 0.04$ & $17.42 \pm 0.03$ & $19.61 \pm 0.04$ \\
		50000 & $3.27 \pm 0.02$ & $12.92 \pm 0.08$ & $17.18 \pm 0.03$ & $19.50 \pm 0.05$ \\
		100000 & $3.24 \pm 0.02$ & $12.07 \pm 0.10$ & $16.66 \pm 0.03$ & $18.60 \pm 0.08$ \\
		\bottomrule
	\end{tabularx}
\end{table}

Графики результатов представлены на рисунках~\ref{img:latency_multi} и~\ref{img:throughput_multi}.

\includeimage
	{latency_multi}
	{f}
	{H}
	{\textwidth}
	{Задержка операций в многопоточном режиме}

\includeimage
	{throughput_multi}
	{f}
	{H}
	{\textwidth}
	{Пропускная способность в многопоточном режиме}

В многопоточном режиме \texttt{sync.Map} показывает лучшую пропускную способность для \texttt{put} (в 2--4 раза), что объясняется оптимизацией \texttt{sync.Map} для сценариев с преобладанием чтения: записи накапливаются в dirty-карте и амортизированно переносятся в read-only кэш.
Для \texttt{get} результаты сравнимы: на малых $N$ разница минимальна, на больших~---~\texttt{sync.Map} имеет небольшое преимущество (10--15\%).

Задержка \texttt{put} в реализации значительно растёт с увеличением $N$ (от 145 до 309 нс/оп), что связано с конкуренцией за \texttt{resizeMu} при частых ресайзах и увеличением длины цепочек в бакетах.
Задержка \texttt{get} растёт значительно медленнее (от 39 до 60 нс/оп), так как чтение не требует захвата \texttt{resizeMu} и использует \texttt{RLock} на бакете.

\section{Потребление памяти}

Таблица~\ref{tab:memory} содержит результаты замеров потребления кучи на операцию для \texttt{put} (для \texttt{get} потребление кучи составляет 0 байт/оп, так как временные аллокации успевают собрать GC между замерами \texttt{HeapAlloc}).

\begin{table}[H]
	\centering
	\caption{Потребление кучи на операцию вставки (байт/оп)}
	\label{tab:memory}
	\begin{tabularx}{\textwidth}{|r|X|X|X|X|}
		\toprule
		$N$ & один поток (concurrent) & один поток (baseline) & много потоков (concurrent) & много потоков (syncmap) \\
		\midrule
		1000 & 311.6 & 160.3 & 313.3 & 138.3 \\
		2500 & 253.3 & 129.7 & 254.1 & 133.8 \\
		5000 & 250.9 & 130.5 & 251.3 & 140.3 \\
		7500 & 324.1 & 167.9 & 324.5 & 142.8 \\
		10000 & 255.2 & 130.8 & 255.5 & 142.2 \\
		25000 & 257.2 & 57.3 & 258.5 & 133.3 \\
		50000 & 260.9 & 84.1 & 260.4 & 135.2 \\
		100000 & 259.1 & 62.6 & 259.3 & 141.9 \\
		\bottomrule
	\end{tabularx}
\end{table}

Графики потребления памяти представлены на рисунках~\ref{img:memory_single} и~\ref{img:memory_multi}.

\includeimage
	{memory_single}
	{f}
	{H}
	{\textwidth}
	{Потребление памяти на операцию (однопоточный режим)}

\includeimage
	{memory_multi}
	{f}
	{H}
	{\textwidth}
	{Потребление памяти на операцию (многопоточный режим)}

Потребление памяти на операцию \texttt{put} в нашей реализации стабильно и составляет около 250--324 байт/оп, что примерно в 2 раза больше, чем у \texttt{map[string]string]} (58--168 байт/оп) и \texttt{sync.Map} (133--142 байт/оп).
Это объясняется оверхедом на \texttt{sync.RWMutex} в каждом бакете (24 байта), структурой \texttt{entry} с двумя строковыми заголовками (32 байта), а также дополнительными аллокациями при росте слайса \texttt{entries} в бакете.
Значение 324 байт/оп при $N = 7500$ объясняется ресайзом таблицы (удвоение числа бакетов с 16 до 32 и перехэширование).

\section{Профилирование}

На рисунках~\ref{img:put_profile} и~\ref{img:get_profile} представлены графы вызовов CPU-профиля для операций \texttt{put} и \texttt{get} соответственно при $N = 50000$, 100 раундов, многопоточный режим.

\includeimage
	{put_profile}
	{f}
	{H}
	{\textwidth}
	{CPU-профиль операции put}

\includeimage
	{get_profile}
	{f}
	{H}
	{\textwidth}
	{CPU-профиль операции get}

В профиле \texttt{put} основное время тратится на:
\begin{itemize}
	\item \texttt{Map.Put}~---~захват мьютексов и поиск/вставка в бакет;
	\item \texttt{Map.maybeResize}~---~проверка и выполнение ресайза;
	\item \texttt{hashIndex}~---~вычисление хэша FNV-1a;
	\item \texttt{memequal}~---~сравнение строковых ключей при поиске в бакете.
\end{itemize}

В профиле \texttt{get} основное время приходится на:
\begin{itemize}
	\item \texttt{Map.Get}~---~захват \texttt{RLock} и линейный поиск в бакете;
	\item \texttt{hashIndex}~---~вычисление хэша;
	\item \texttt{memequal}~---~сравнение ключей.
\end{itemize}

\chapter{Вывод}

Реализована конкурентная хэш-таблица на основе striped locking с \texttt{sync.RWMutex} на каждый бакет, \texttt{atomic.Value} для lock-free доступа к таблице и динамическим ресайзом при превышении порога загрузки.

Проведены замеры производительности на наборах данных от 1000 до 100000 ключей в однопоточном и многопоточном режимах.
Можно сделать следующие выводы:
\begin{itemize}
	\item в однопоточном режиме реализация в 2--4 раза медленнее \texttt{map[string]string} для \texttt{put} и в 1.4--2.8 раза для \texttt{get} из-за накладных расходов на синхронизацию;
	\item в многопоточном режиме \texttt{sync.Map} показывает лучшую пропускную способность для \texttt{put} (в 2--4 раза), тогда как для \texttt{get} результаты сравнимы (отставание 10--15\%);
	\item потребление памяти на операцию \texttt{put} стабильно и составляет около 250--324 байт/оп, что примерно в 2 раза больше, чем у \texttt{map[string]string} и \texttt{sync.Map};
	\item профилирование подтверждает, что основное время тратится на захват мьютексов, вычисление хэша и сравнение ключей в бакетах.
\end{itemize}

Основным узким местом реализации является глобальная блокировка \texttt{resizeMu} при ресайзе, которая блокирует все операции \texttt{Put}.
Для улучшения масштабируемости можно применить переносимые блокировки (lock striping с переносом) или lock-free схему ресайза.

\end{document}
